\documentclass[11pt, a4paper]{article}

% ── Packages ──────────────────────────────────────────────────
\usepackage[utf8]{inputenc}
\usepackage[T1]{fontenc}
\usepackage{lmodern}
\usepackage[margin=2.5cm]{geometry}
\usepackage{amsmath, amssymb, amsthm, mathtools}
\usepackage{booktabs}
\usepackage{hyperref}
\usepackage{enumitem}
\usepackage{fancyhdr}
\usepackage[most]{tcolorbox}
\usepackage{xcolor}
\usepackage{listings}
\usepackage{array}

% ── Colors ────────────────────────────────────────────────────
\definecolor{ddpmblue}{RGB}{40, 80, 160}
\definecolor{fmgreen}{RGB}{30, 120, 60}
\definecolor{noteorange}{RGB}{200, 100, 0}
\definecolor{codebg}{RGB}{245, 245, 245}
\definecolor{codeframe}{RGB}{200, 200, 200}

% ── Boxes ─────────────────────────────────────────────────────
\newtcolorbox{keyinsight}{
  colback=fmgreen!8, colframe=fmgreen!70!black,
  fonttitle=\bfseries, title=Key Insight,
  rounded corners, boxrule=0.8pt
}
\newtcolorbox{ddpmbox}{
  colback=ddpmblue!6, colframe=ddpmblue!70!black,
  fonttitle=\bfseries, title=DDPM,
  rounded corners, boxrule=0.8pt
}
\newtcolorbox{fmbox}{
  colback=fmgreen!6, colframe=fmgreen!70!black,
  fonttitle=\bfseries, title=Flow Matching,
  rounded corners, boxrule=0.8pt
}
\newtcolorbox{warningbox}{
  colback=noteorange!8, colframe=noteorange!70!black,
  fonttitle=\bfseries, title=Watch Out,
  rounded corners, boxrule=0.8pt
}

% ── Code listings ─────────────────────────────────────────────
\lstset{
  language=Python,
  basicstyle=\ttfamily\small,
  backgroundcolor=\color{codebg},
  frame=single,
  rulecolor=\color{codeframe},
  keywordstyle=\color{ddpmblue}\bfseries,
  commentstyle=\color{gray},
  stringstyle=\color{fmgreen},
  breaklines=true,
  columns=fullflexible,
  xleftmargin=1em,
  framexleftmargin=0.5em,
  aboveskip=1em,
  belowskip=1em,
}

% ── Math shortcuts ────────────────────────────────────────────
\newcommand{\vtheta}{v_\theta}
\newcommand{\mutheta}{\mu_\theta}
\newcommand{\mutilde}{\tilde{\mu}}
\newcommand{\betatilde}{\tilde{\beta}}
\newcommand{\alphabar}{\bar{\alpha}}
\newcommand{\R}{\mathbb{R}}
\newcommand{\E}{\mathbb{E}}
\newcommand{\N}{\mathcal{N}}
\newcommand{\Id}{\mathbf{I}}
\newcommand{\norm}[1]{\left\lVert #1 \right\rVert}
\DeclareMathOperator{\Div}{div}
\DeclareMathOperator{\MSE}{MSE}

% ── Header/Footer ─────────────────────────────────────────────
\pagestyle{fancy}
\fancyhf{}
\fancyhead[L]{\small Flow Matching Reading Guide}
\fancyhead[R]{\small ToyDiffusion}
\fancyfoot[C]{\thepage}

% ── Hyperlinks ────────────────────────────────────────────────
\hypersetup{
  colorlinks=true,
  linkcolor=ddpmblue,
  urlcolor=fmgreen!80!black,
  citecolor=ddpmblue,
  pdfstartview=FitH,
}

% ══════════════════════════════════════════════════════════════
\title{
  \textbf{Flow Matching: Complete Reading Guide} \\[0.3em]
  \large For Upgrading ToyDiffusion from DDPM to Flow Matching
}
\author{Compiled for offline study}
\date{February 2026}

\begin{document}
\maketitle
\thispagestyle{fancy}

\begin{abstract}
This document is a self-contained reference for understanding how to upgrade
ToyDiffusion---a toy DDPM implementation that generates 2D Homer Simpson
scatter plots---from the classic \textbf{DDPM} framework (Ho et al., 2020) to
the modern \textbf{Flow Matching} paradigm (Lipman et al., 2022; Liu et al.,
2022). All key equations are typeset in full, derivations are expanded, and a
side-by-side comparison makes the correspondence explicit. Stable Diffusion~3,
Flux, and most modern production generative models have switched from DDPM to
flow matching.
\end{abstract}

\tableofcontents
\newpage

% ══════════════════════════════════════════════════════════════
\section{The Big Picture}
\label{sec:big-picture}

\subsection{What DDPM Does (Your Original Code)}

Your ToyDiffusion project implements \textbf{DDPM} (Ho et al., 2020):

\begin{enumerate}
  \item Define a \textbf{Markov chain} that adds Gaussian noise over $T$
        discrete steps.
  \item Train a neural network to \textbf{reverse} the process step by step.
  \item The model predicts either the noise $\epsilon$ or the posterior mean
        $\mutilde_t$.
\end{enumerate}

The forward process gradually destroys the data signal:
\[
  q(x_t \mid x_0)
  = \N\!\bigl(x_t;\; \sqrt{\alphabar_t}\, x_0,\;
    (1-\alphabar_t)\,\Id\bigr)
\]
and the reverse process learns to undo it, one step at a time.

\subsection{What Flow Matching Does}

\textbf{Flow Matching} (2022) provides a simpler alternative:

\begin{enumerate}
  \item Define a \textbf{continuous path} from noise to data via linear
        interpolation.
  \item Train a neural network to predict the \textbf{velocity field}
        $\vtheta(x,t)$ along that path.
  \item Generate samples by integrating a simple ODE.
\end{enumerate}

\begin{keyinsight}
Diffusion models are a \emph{special case} of flow matching with a particular
(suboptimal) choice of interpolation path.  Using straight-line (Optimal
Transport) paths instead gives straighter trajectories, which means fewer
sampling steps and simpler math.
\end{keyinsight}

% ══════════════════════════════════════════════════════════════
\section{Paper 1: Flow Matching for Generative Modeling}
\label{sec:paper1}

\begin{quote}
\textbf{Authors:} Yaron Lipman, Ricky T.~Q.~Chen, Heli Ben-Hamu,
Maximilian Nickel, Matt Le (Meta FAIR) \\
\textbf{Published:} October 2022 (ICLR 2023) \\
\textbf{arXiv:} \url{https://arxiv.org/abs/2210.02747}
\end{quote}

\subsection{Density Functions vs.\ Samples: A Key Distinction}
\label{sec:density-vs-samples}

Before diving into the math, we need to be precise about one thing.
A \textbf{probability density function} $p(x)$ is a mathematical function
$p: \R^d \to \R$ satisfying $\int p(x)\,dx = 1$.  It tells you ``how likely
is $x$?''  A \textbf{sample} (or realization) is an actual point
$x \in \R^d$ drawn from that distribution.

In your ToyDiffusion project:

\begin{itemize}
  \item \textbf{The Gaussian prior} $p_0 = \N(0, \Id)$:
    we have the exact formula
    $p_0(x) = \frac{1}{(2\pi)^{d/2}}\exp\!\bigl(-\tfrac{1}{2}\|x\|^2\bigr)$,
    \emph{and} we can generate samples via \texttt{torch.randn(N, 2)}.

  \item \textbf{Homer Simpson} $p_{\mathrm{data}}$:
    we do \textbf{not} have a formula.  We only have samples---the $(x,y)$
    scatter plot points from your dataset.  But mathematically, some unknown
    density function $p_{\mathrm{data}}(x)$ exists that is high where Homer's
    features are and low everywhere else.

  \item \textbf{Intermediate distributions} $p_t$ (for $0 < t < 1$):
    we have \textbf{neither} the formula \textbf{nor} pre-computed samples.
    We \emph{could} generate samples at time $t$ by flowing $p_0$ samples
    forward, but we never write down $p_t(x)$ as a closed-form expression.
\end{itemize}

\begin{warningbox}
The math in this document talks about density functions $p_t(x)$, but your
code only ever works with \textbf{samples} (actual tensors of points).
The entire point of Conditional Flow Matching (Section~\ref{sec:cfm}) is to
set up a training loss that only requires samples from $p_0$ and
$p_{\mathrm{data}}$---never $p_t$ itself.
\end{warningbox}


\subsection{Continuous Normalizing Flows}

A \textbf{Continuous Normalizing Flow} (CNF) defines a time-dependent ODE
whose solution transforms a simple prior $p_0 = \N(0,\Id)$ into the data
distribution $p_1 \approx p_{\mathrm{data}}$:
\begin{equation}
  \label{eq:ode}
  \frac{dx}{dt} = \vtheta(x, t), \qquad t \in [0, 1].
\end{equation}

\paragraph{What ``integrating this ODE'' means.}
Recall that solving any ODE requires two things:
\begin{enumerate}
  \item The differential equation itself: $\frac{dx}{dt} = \vtheta(x, t)$
  \item An \textbf{initial condition}: $x(0) = x_0$
\end{enumerate}

Given an initial sample $x_0 \sim \N(0, \Id)$, the solution at time $t{=}1$
is:
\begin{equation}
  \label{eq:ode-solution}
  x_1 = x_0 + \int_0^1 \vtheta\!\bigl(x(t),\, t\bigr)\, dt
\end{equation}

Note: $\int_0^1 \frac{dx}{dt}\,dt = x(1) - x(0) = x_1 - x_0$ gives you the
\emph{change} in $x$, not $x_1$ itself.  You always need the initial condition
$x_0$ to recover $x_1$.  In practice this integral is approximated by Euler
steps (see Section~\ref{sec:code-changes}):
\begin{equation}
  x_{t+\Delta t} \approx x_t + \vtheta(x_t, t) \cdot \Delta t
\end{equation}

\paragraph{The flow map.}
The ODE defines a \textbf{flow map} $\phi_t : \R^d \to \R^d$---a function
that tells you where each point ends up at time~$t$:
\[
  \phi_0(x) = x \quad \text{(identity)}, \qquad
  \frac{d}{dt}\phi_t(x) = \vtheta\!\bigl(\phi_t(x),\, t\bigr).
\]
If you start at $x_0$, then $\phi_t(x_0)$ is your position at time $t$.

\paragraph{The pushforward (what $[\phi_t]_\# \, p_0$ means).}
This notation looks scary but the idea is simple: take every point sampled
from $p_0$, apply the map $\phi_t$ to each point, and look at the
distribution of the results.  That resulting distribution is the
\textbf{pushforward}:
\begin{equation}
  p_t = [\phi_t]_\# \, p_0.
\end{equation}

In code, this is exactly what your ToyDiffusion visualization already does:
\begin{lstlisting}
# Sample 1000 points from p_0
x = torch.randn(1000, 2)          # these are distributed as p_0

# Flow them forward to time t using the ODE
for step in range(n_steps_to_t):
    x = x + v_theta(x, t) * dt    # Euler integration

# Now x is distributed as p_t = [phi_t]_# p_0
plot(x)  # visualize the pushforward!
\end{lstlisting}

When you animated particles flowing from noise to Homer Simpson, you were
computing the pushforward at each frame.


\subsection{The Continuity Equation}
\label{sec:continuity}

The probability path $p_t$ and velocity field $u_t$ must satisfy the
\textbf{continuity equation}.  Before stating it, let's unpack the notation.

\paragraph{Notation: $u_t$ is the velocity field.}
The papers use $u_t(x)$ as their notation for the velocity field.  This is
the same object as $\vtheta(x,t)$ (what your code will call it)---just
different notation.  Don't let the symbol change confuse you:
$u_t(x) = \vtheta(x,t) = v(x,t)$ all refer to the same thing.

\paragraph{What is divergence?}
The divergence $\Div(\cdot)$ measures \emph{how much a vector field is
spreading out or converging} at a point.  For a 2D velocity field
$v = (v_x, v_y)$:
\[
  \Div(v) = \frac{\partial v_x}{\partial x} + \frac{\partial v_y}{\partial y}
\]

\begin{itemize}
  \item $\Div > 0$: particles are spreading apart (like a source / explosion)
  \item $\Div < 0$: particles are converging (like a sink / drain)
  \item $\Div = 0$: the flow preserves volume (incompressible)
\end{itemize}

\paragraph{The equation.}
Now we can state it:
\begin{equation}
  \label{eq:continuity}
  \underbrace{\frac{\partial p_t}{\partial t}}_{\substack{
    \text{how fast density} \\ \text{changes at } x}}
  + \underbrace{\Div\bigl(p_t \cdot u_t\bigr)}_{\substack{
    \text{how much probability} \\ \text{mass flows away from } x}}
  = 0.
\end{equation}

The product $p_t \cdot u_t$ is the \textbf{probability flux}: density times
velocity = ``flow of probability mass.''  The equation says these two terms
must balance: if probability mass is flowing away from a region
($\Div(p_t \cdot u_t) > 0$), the density there must decrease
($\frac{\partial p_t}{\partial t} < 0$), and vice versa.

\begin{keyinsight}
This is just the law of conservation of probability mass.
Probability can't be created or destroyed, only moved around.
If particles leave a region, density decreases there.
If particles enter, density increases.
The continuity equation is the mathematical statement of this obvious
physical fact.
\end{keyinsight}

\subsection{The Flow Matching Loss (Intractable)}

The ideal objective would be to match the velocity field directly:
\begin{equation}
  \label{eq:lfm}
  \mathcal{L}_{\mathrm{FM}}(\theta)
  = \E_{t \sim U(0,1),\; x \sim p_t}
    \norm{\vtheta(x, t) - u_t(x)}^2
\end{equation}
where $u_t(x)$ is the true velocity field that generates the probability path
$p_t$.

\begin{warningbox}
This loss is \textbf{intractable}: we don't have access to $p_t(x)$ or
$u_t(x)$ in closed form, because they involve integrals over all data points.
Remember from Section~\ref{sec:density-vs-samples}: $p_t$ is a theoretical
density we never compute!  We can't sample from it, and we don't know $u_t$.
This is why we need the conditional trick below.
\end{warningbox}

\subsection{Conditional Flow Matching (The Key Contribution)}
\label{sec:cfm}

Instead of working with intractable marginals, we define \textbf{conditional}
probability paths.  For each data point $x_1 \sim p_{\mathrm{data}}$, the
conditional probability path is:
\begin{equation}
  \label{eq:cond-path}
  p_t(x \mid x_1) = \N\!\bigl(x;\; \mu_t(x_1),\; \sigma_t(x_1)^2 \,\Id\bigr)
\end{equation}
with boundary conditions:
\begin{align}
  t = 0: &\quad \mu_0 = 0,\quad \sigma_0 = 1
         \quad \text{(standard Gaussian prior)} \\
  t = 1: &\quad \mu_1 = x_1,\quad \sigma_1 \approx 0
         \quad \text{(concentrated at data point)}
\end{align}

For the \textbf{Optimal Transport (OT) linear path}:
\begin{equation}
  \label{eq:ot-path}
  \mu_t(x_1) = t \cdot x_1, \qquad \sigma_t = 1 - t
\end{equation}

Sampling from this conditional path is equivalent to:
\begin{equation}
  \label{eq:interpolation}
  \boxed{x_t = (1 - t)\, x_0 + t\, x_1}
  \qquad \text{where } x_0 \sim \N(0, \Id)
\end{equation}

The conditional velocity field for this path is:
\begin{equation}
  \label{eq:cond-velocity-general}
  u_t(x \mid x_1) = \frac{\sigma_t'(x_1)}{\sigma_t(x_1)}
    \bigl(x - \mu_t(x_1)\bigr) + \mu_t'(x_1)
\end{equation}

\paragraph{Derivation.}
Taking the time derivative of $x_t = (1-t)\,x_0 + t\,x_1$:
\begin{equation}
  \frac{dx_t}{dt} = -x_0 + x_1 = x_1 - x_0
\end{equation}

So the conditional velocity for the OT path simplifies to the remarkably simple
expression:
\begin{equation}
  \label{eq:cond-velocity-ot}
  \boxed{u_t(x \mid x_1) = x_1 - x_0}
\end{equation}

A constant!  The velocity doesn't depend on $x$ or $t$---it's simply the
straight-line direction from noise to data.

The \textbf{Conditional Flow Matching (CFM)} loss is then:
\begin{equation}
  \label{eq:lcfm}
  \boxed{
    \mathcal{L}_{\mathrm{CFM}}(\theta)
    = \E_{t \sim U(0,1),\; x_0 \sim \N(0,\Id),\; x_1 \sim p_{\mathrm{data}}}
      \norm{\vtheta(x_t, t) - (x_1 - x_0)}^2
  }
\end{equation}

\begin{keyinsight}
\textbf{Key Theorem (Lipman et al., 2022):}
The gradients of $\mathcal{L}_{\mathrm{CFM}}$ and $\mathcal{L}_{\mathrm{FM}}$
are identical:
\[
  \nabla_\theta \,\mathcal{L}_{\mathrm{CFM}}(\theta)
  = \nabla_\theta \,\mathcal{L}_{\mathrm{FM}}(\theta).
\]
Training with the tractable conditional loss is \emph{equivalent} to training
with the intractable marginal loss.  This is the central result that makes flow
matching practical.
\end{keyinsight}

\paragraph{Proof sketch.}
The marginal velocity field decomposes as:
\[
  u_t(x) = \frac{\int u_t(x \mid x_1)\, p_t(x \mid x_1)\, q(x_1)\, dx_1}
                 {\int p_t(x \mid x_1)\, q(x_1)\, dx_1}
         = \E_{q(x_1 \mid x)}\bigl[u_t(x \mid x_1)\bigr]
\]
Expanding the FM loss and applying the tower property of conditional
expectations shows the gradient equivalence.  The full proof relies on
exchanging the order of integration over $x$ and $x_1$.

\subsection{Why Optimal Transport Paths?}

The linear interpolation $x_t = (1-t)\,x_0 + t\,x_1$ is the
\textbf{Wasserstein-2 optimal transport geodesic} between the conditional
distributions $\N(0,\Id)$ and $\delta_{x_1}$.

This has two important consequences:
\begin{enumerate}
  \item The marginal velocity field is smoother and closer to straight.
  \item Fewer ODE solver steps are needed at inference (even 10--20 Euler steps
        work well).
\end{enumerate}

Compare: DDPM-style paths are curved (they first destroy the signal heavily,
then denoise), requiring ${\sim}1000$ steps or specialized samplers like DDIM.

\subsection{Connection to Diffusion Models}

Diffusion models correspond to flow matching with a \textbf{specific non-linear
path}:
\begin{equation}
  \label{eq:vp-path}
  x_t = \sqrt{\alphabar_t}\, x_1 + \sqrt{1 - \alphabar_t}\;\epsilon,
  \qquad \epsilon \sim \N(0, \Id)
\end{equation}

This is the Variance Preserving SDE (VP-SDE) schedule.  In the flow matching
framework, this corresponds to:
\[
  \mu_t(x_1) = \sqrt{\alphabar_t}\, x_1,
  \qquad
  \sigma_t = \sqrt{1 - \alphabar_t}
\]

The score function and velocity field are related by:
\begin{equation}
  \label{eq:score-velocity}
  v(x, t) = f(x, t) - \tfrac{1}{2}\, g(t)^2 \,\nabla_x \log p_t(x)
\end{equation}
where $f$ and $g$ are the drift and diffusion coefficients of the corresponding
SDE formulation $dx = f(x,t)\,dt + g(t)\,dW_t$.

% ══════════════════════════════════════════════════════════════
\section{Paper 2: Rectified Flow}
\label{sec:paper2}

\begin{quote}
\textbf{Authors:} Xingchao Liu, Chengyue Gong, Qiang Liu (UT Austin) \\
\textbf{Published:} September 2022 (ICLR 2023 Spotlight) \\
\textbf{arXiv:} \url{https://arxiv.org/abs/2209.03003}
\end{quote}

\subsection{Core Formulation}

Same ODE setup: learn $\vtheta$ such that
\begin{equation}
  \frac{dZ_t}{dt} = \vtheta(Z_t, t), \qquad t \in [0, 1].
\end{equation}

The training objective is identical to Conditional Flow Matching:
\begin{equation}
  \label{eq:rectified-loss}
  \mathcal{L}
  = \E_{t \sim U(0,1)}\;
    \E_{X_0 \sim \N(0,\Id),\; X_1 \sim p_{\mathrm{data}}}
    \norm{\vtheta(X_t, t) - (X_1 - X_0)}^2
\end{equation}
where $X_t = (1 - t)\, X_0 + t\, X_1$.

\subsection{The Reflow Procedure (Unique Contribution)}

The key idea: \textbf{iteratively straighten the flow trajectories}.

\paragraph{Round 1 (Initial Training).}
Sample independent pairs $(X_0, X_1)$ where $X_0 \sim \N(0,\Id)$ and
$X_1 \sim p_{\mathrm{data}}$.  Train $v_1$ with the standard loss
\eqref{eq:rectified-loss}.

\paragraph{Round 2 (Reflow).}
For many samples $X_0 \sim \N(0,\Id)$, solve the ODE of $v_1$ forward to
obtain:
\begin{equation}
  Z_1 = X_0 + \int_0^1 v_1(Z_t, t)\, dt
\end{equation}

Now $(X_0, Z_1)$ are \textbf{coupled} pairs---they lie on the same learned
trajectory.  Retrain a new model $v_2$ on these coupled pairs:
\begin{equation}
  \mathcal{L}_{\mathrm{reflow}}
  = \E_{t}\, \norm{v_2\bigl((1{-}t)\,X_0 + t\,Z_1,\; t\bigr)
    - (Z_1 - X_0)}^2
\end{equation}

\paragraph{Why it works.}
In round~1, independent pairs produce crossing linear interpolation paths,
forcing the velocity field to average at intersections (creating curvature).
After reflow, the pairs are correlated---their interpolations cross much less,
so the retrained flow is straighter.

\subsection{Theoretical Guarantees}

The straightness of the flow is measured by:
\begin{equation}
  \label{eq:straightness}
  S(v) = \int_0^1 \E\!\left[
    \norm{v(Z_t, t) - (Z_1 - Z_0)}^2
  \right] dt
\end{equation}

\begin{itemize}
  \item $S(v) = 0$ if and only if all trajectories are perfectly straight
        lines.
  \item Each round of reflow \textbf{monotonically decreases} $S(v)$, i.e.\
        $S(v_{k+1}) \le S(v_k)$.
  \item In the limit, $S(v) \to 0$ and all trajectories become straight.
\end{itemize}

\subsection{Few-Step and One-Step Generation}

If all trajectories are perfectly straight, a single Euler step suffices:
\begin{equation}
  \label{eq:one-step}
  Z_1 = Z_0 + v(Z_0, 0)
  \qquad \text{(one-step generation!)}
\end{equation}

In practice:
\begin{itemize}
  \item After 1 reflow: good results with 2--4 Euler steps.
  \item After 2 reflows: competitive 1-step generation.
\end{itemize}

% ══════════════════════════════════════════════════════════════
\section{DDPM vs.\ Flow Matching: Side-by-Side}
\label{sec:comparison}

\subsection{Notation Comparison}

\begin{warningbox}
The time direction is \textbf{flipped}!  In DDPM, $t{=}0$ is clean data and
$t{=}T$ is noise.  In flow matching, $t{=}0$ is noise and $t{=}1$ is data.
\end{warningbox}

\begin{center}
\renewcommand{\arraystretch}{1.3}
\begin{tabular}{@{} l c c @{}}
\toprule
\textbf{Concept} & \textbf{DDPM (Your Code)} & \textbf{Flow Matching} \\
\midrule
Time domain & $t \in \{0, 1, \ldots, T\}$ discrete & $t \in [0, 1]$ continuous \\
Data & $x_0$ (original, clean) & $x_1$ (target at $t{=}1$) \\
Noise & $x_T$ (fully noised) & $x_0$ (noise at $t{=}0$) \\
Direction & noise $\to$ data: $t\!:\, T \to 0$ & noise $\to$ data: $t\!:\, 0 \to 1$ \\
\bottomrule
\end{tabular}
\end{center}

\subsection{Process Comparison}

\begin{center}
\renewcommand{\arraystretch}{1.4}
\begin{tabular}{@{} >{\raggedright\arraybackslash}p{3.2cm}
                    >{\raggedright\arraybackslash}p{5.5cm}
                    >{\raggedright\arraybackslash}p{5.5cm} @{}}
\toprule
\textbf{Aspect} & \textbf{DDPM (Your Original)} & \textbf{Flow Matching} \\
\midrule
Forward process &
  Markov chain, $T$ discrete steps &
  Linear interpolation, continuous \\
Noise schedule &
  Required ($\beta_t$, $\alpha_t$, $\alphabar_t$) &
  \textbf{Not needed} \\
Training target &
  Predict $\epsilon$ or $\mutilde_t$ &
  Predict velocity $v = x_1 - x_0$ \\
Loss function &
  $\norm{\mutilde_t - \mutheta}^2$ &
  $\norm{v_{\mathrm{target}} - \vtheta}^2$ \\
Sampling &
  ${\sim}1000$ steps, complex posterior &
  ${\sim}10$--$100$ Euler steps: $x \mathrel{{+}{=}} v \cdot dt$ \\
Math complexity &
  KL divergence, posterior distributions &
  Simple ODE integration \\
\bottomrule
\end{tabular}
\end{center}

\subsection{Equation Comparison}

\subsubsection*{Interpolation (creating a noisy sample at time~$t$)}

\begin{ddpmbox}
\vspace{-0.5em}
\[
  x_t = \sqrt{\alphabar_t}\, x_0
      + \sqrt{1 - \alphabar_t}\;\epsilon,
  \qquad \epsilon \sim \N(0, \Id)
\]
\end{ddpmbox}

\begin{fmbox}
\vspace{-0.5em}
\[
  x_t = (1 - t)\, x_0 + t\, x_1,
  \qquad x_0 \sim \N(0, \Id)
\]
\end{fmbox}

\subsubsection*{Training target (what the network predicts)}

\begin{ddpmbox}
\vspace{-0.3em}
\textbf{$\mu$-prediction:} a complex function of $x_0$, $x_t$, $\alpha_t$, $\beta_t$:
\begin{align*}
  \mutilde_t
  &= \frac{\sqrt{\alphabar_{t-1}}\;\beta_t}{1 - \alphabar_t}\,x_0
  + \frac{\sqrt{\alpha_t}\;(1 - \alphabar_{t-1})}{1 - \alphabar_t}\,x_t
\end{align*}

\textbf{$\epsilon$-prediction:} the noise that was added: target $= \epsilon$.
\end{ddpmbox}

\begin{fmbox}
\vspace{-0.3em}
\textbf{Velocity prediction:} the direction from noise to data:
\[
  v_{\mathrm{target}} = x_1 - x_0
\]
\end{fmbox}

\subsubsection*{Sampling (generating new data)}

\begin{ddpmbox}
\vspace{-0.5em}
Reverse Markov chain:
\[
  x_{t-1} = \mutheta(x_t, t) + \sigma_t \, z,
  \qquad z \sim \N(0, \Id)
\]
\end{ddpmbox}

\begin{fmbox}
\vspace{-0.5em}
Euler ODE step:
\[
  x_{t + \Delta t} = x_t + \vtheta(x_t, t) \cdot \Delta t
\]
\end{fmbox}

% ══════════════════════════════════════════════════════════════
\section{The Math You Need}
\label{sec:math}

\subsection{Continuous Normalizing Flows and the Continuity Equation}

A CNF defines a velocity field $v: \R^d \times [0,1] \to \R^d$ whose ODE
\eqref{eq:ode} has the unique solution flow $\phi_t$.  The induced probability
path $p_t$ satisfies the continuity equation~\eqref{eq:continuity}.

The \textbf{change of variables formula} for CNFs gives the exact log-likelihood:
\begin{equation}
  \log p_1(x_1) = \log p_0\bigl(\phi_1^{-1}(x_1)\bigr)
  - \int_0^1 \Div\bigl(v(\phi_t(x), t)\bigr)\, dt
\end{equation}

This is important theoretically, though in practice flow matching sidesteps
likelihood computation by using the regression loss \eqref{eq:lcfm} directly.

\subsection{Gaussian Conditional Probability Paths}

The general Gaussian conditional path is:
\begin{equation}
  p_t(x \mid x_1) = \N\!\bigl(x;\; \mu_t(x_1),\; \sigma_t(x_1)^2\,\Id\bigr)
\end{equation}

with boundary conditions:
\begin{alignat}{3}
  &t = 0: &\qquad \mu_0 &= 0, &\quad \sigma_0 &= 1
    \quad \text{(standard Gaussian prior)} \\
  &t = 1: &\qquad \mu_1 &= x_1, &\quad \sigma_1 &\approx 0
    \quad \text{(concentrated at data)}
\end{alignat}

Two important choices:

\renewcommand{\arraystretch}{1.5}
\begin{center}
\begin{tabular}{@{} l l l @{}}
\toprule
\textbf{Path} & $\mu_t(x_1)$ & $\sigma_t$ \\
\midrule
OT / Linear (recommended) & $t \cdot x_1$ & $1 - t$ \\
VP / Diffusion (DDPM-like) & $\sqrt{\alphabar_t}\, x_1$
  & $\sqrt{1 - \alphabar_t}$ \\
\bottomrule
\end{tabular}
\end{center}

Sampling from either path follows:
\begin{equation}
  x_t = \mu_t(x_1) + \sigma_t \cdot \epsilon, \qquad \epsilon \sim \N(0, \Id)
\end{equation}

For the OT path: $x_t = t\, x_1 + (1-t)\,\epsilon = (1-t)\,x_0 + t\,x_1$
(with $x_0 = \epsilon$).

\subsection{Deriving the Conditional Velocity Field}

\paragraph{General case.}  Given the Gaussian conditional path
$p_t(x \mid x_1) = \N(x;\, \mu_t, \sigma_t^2\, \Id)$, we can write
$x_t = \mu_t + \sigma_t\, \epsilon$ where $\epsilon \sim \N(0,\Id)$ is
fixed.  Taking the time derivative:
\begin{equation}
  u_t(x \mid x_1) = \frac{dx_t}{dt}
  = \mu_t' + \sigma_t'\, \epsilon
  = \mu_t' + \frac{\sigma_t'}{\sigma_t}\,(x_t - \mu_t)
\end{equation}

where in the last step we used $\epsilon = (x_t - \mu_t)/\sigma_t$.

\paragraph{OT/Linear path.}
With $\mu_t = t\,x_1$ and $\sigma_t = 1-t$, we get
$\mu_t' = x_1$ and $\sigma_t' = -1$, so:
\begin{align}
  u_t(x \mid x_1)
  &= x_1 + \frac{-1}{1-t}\,(x_t - t\, x_1) \notag \\
  &= x_1 - \frac{x_t - t\, x_1}{1-t} \notag \\
  &= x_1 - \frac{(1-t)\,x_0 + t\,x_1 - t\,x_1}{1-t}
    \quad \text{(substituting $x_t = (1-t)x_0 + t\,x_1$)} \notag \\
  &= x_1 - x_0
\end{align}

\paragraph{VP/Diffusion path.}
With $\mu_t = \sqrt{\alphabar_t}\, x_1$ and
$\sigma_t = \sqrt{1-\alphabar_t}$:
\begin{equation}
  u_t^{\mathrm{VP}}(x \mid x_1)
  = \frac{d\sqrt{\alphabar_t}}{dt}\, x_1
  + \frac{d\sqrt{1-\alphabar_t}/dt}{\sqrt{1-\alphabar_t}}
    \,(x - \sqrt{\alphabar_t}\, x_1)
\end{equation}

This is significantly more complex, depending on the specific $\alphabar_t$
schedule.

\subsection{Score--Velocity--Epsilon Conversion}

Given a general schedule $\alpha_t$ (not to be confused with $\alphabar_t$
in DDPM notation), where $x_t = \alpha_t\, x_1 + \sigma_t\, \epsilon$, the
three prediction targets are related:

\begin{align}
  \label{eq:conversions}
  \text{Velocity:} \quad
    v &= \alpha_t'\, x_1 + \sigma_t'\, \epsilon \\[0.5em]
  \text{Score:} \quad
    \nabla_x \log p_t(x) &= -\frac{x - \alpha_t\, x_1}{\sigma_t^2}
    = -\frac{\epsilon}{\sigma_t} \\[0.5em]
  \text{Epsilon:} \quad
    \epsilon &= \frac{x_t - \alpha_t\, x_1}{\sigma_t} \\[0.5em]
  \text{Data:} \quad
    x_1 &= \frac{x_t - \sigma_t\, \epsilon}{\alpha_t}
\end{align}

You can convert between any two given the schedule $(\alpha_t, \sigma_t)$.

\paragraph{Velocity from epsilon (for OT path, $\alpha_t = t$, $\sigma_t = 1-t$):}
\begin{equation}
  v = \alpha_t'\, x_1 + \sigma_t'\, \epsilon
    = 1 \cdot x_1 + (-1) \cdot \epsilon
    = x_1 - \epsilon = x_1 - x_0
\end{equation}

\paragraph{Velocity from epsilon (for VP path):}
\begin{equation}
  v = \frac{d\sqrt{\alphabar_t}}{dt}\, x_1
    - \frac{d\sqrt{1-\alphabar_t}}{dt}\;\epsilon
\end{equation}

% ══════════════════════════════════════════════════════════════
\section{The 3 Changes to Your Code}
\label{sec:code-changes}

\subsection{Change 1: Replace \texttt{q\_sample} with Linear Interpolation}

\begin{ddpmbox}
\vspace{-0.5em}
Old (\texttt{diffusion.py:q\_sample}):
\begin{lstlisting}
# x_t = sqrt(alpha_bar_t) * x_0 + sqrt(1 - alpha_bar_t) * eps
alpha_bar_t = list_bar_alphas[t]
mean = alpha_bar_t * x_start
cov = torch.eye(x_start.shape[0]) * (1 - alpha_bar_t)
x_t = MultivariateNormal(mean, cov).sample()
\end{lstlisting}
\end{ddpmbox}

\begin{fmbox}
\vspace{-0.5em}
New (flow matching interpolation):
\begin{lstlisting}
def interpolate(x_data, t):
    x_noise = torch.randn_like(x_data)       # x_0 ~ N(0, I)
    x_t = (1 - t) * x_noise + t * x_data     # Eq. (6)
    v_target = x_data - x_noise               # Eq. (9)
    return x_t, v_target
\end{lstlisting}
\end{fmbox}

No more $\alphabar_t$, $\alpha_t$, or $\beta_t$ schedules!

\subsection{Change 2: Training Target $\to$ Velocity}

\begin{ddpmbox}
\vspace{-0.5em}
Old (trains $\mutheta$ against $\mutilde_t$):
\begin{lstlisting}
mu_theta = model(x_t, t)
mu_tilde, _ = posterior_q(x_start, x_t, t, ...)
loss = MSE(mu_tilde, mu_theta)
\end{lstlisting}
\end{ddpmbox}

\begin{fmbox}
\vspace{-0.5em}
New (trains $\vtheta$ against $v_{\mathrm{target}} = x_1 - x_0$):
\begin{lstlisting}
v_pred = model(x_t, t)
loss = MSE(v_target, v_pred)
\end{lstlisting}
\end{fmbox}

\subsection{Change 3: Sampling $\to$ Euler ODE Integration}

\begin{ddpmbox}
\vspace{-0.5em}
Old (complex reverse Markov chain):
\begin{lstlisting}
# For t = T, T-1, ..., 1:
mu_theta = model(x_t, t)
x_{t-1} ~ N(mu_theta, beta_t * I)
\end{lstlisting}
\end{ddpmbox}

\begin{fmbox}
\vspace{-0.5em}
New (simple Euler integration):
\begin{lstlisting}
x = torch.randn(n_samples, dim)   # start from noise
dt = 1.0 / n_steps
for i in range(n_steps):
    t = i * dt
    v = model(x, t)
    x = x + v * dt                # Euler step
# x is now a generated sample
\end{lstlisting}
\end{fmbox}

\subsection{Complete Training Algorithm}

\begin{tcolorbox}[colback=codebg, colframe=codeframe, title=\textbf{Algorithm 1: Flow Matching Training}]
\begin{lstlisting}
for each batch:
    x_1 = sample_data()                # from dataset
    x_0 = torch.randn_like(x_1)        # from N(0, I)
    t = torch.rand(batch_size, 1)       # from U(0, 1)
    x_t = (1 - t) * x_0 + t * x_1      # interpolation
    v_target = x_1 - x_0               # target velocity
    v_pred = model(x_t, t)             # network prediction
    loss = MSE(v_pred, v_target)
    loss.backward()
    optimizer.step()
\end{lstlisting}
\end{tcolorbox}

\subsection{Complete Sampling Algorithm}

\begin{tcolorbox}[colback=codebg, colframe=codeframe, title=\textbf{Algorithm 2: Flow Matching Sampling (Euler)}]
\begin{lstlisting}
x = torch.randn(n_samples, dim)        # start from noise
dt = 1.0 / n_steps
for i in range(n_steps):
    t = i * dt
    v = model(x, t)
    x = x + v * dt                     # Euler step
# x is now a generated sample
\end{lstlisting}
\end{tcolorbox}

% ══════════════════════════════════════════════════════════════
\section{Advanced: Reflow for One-Step Generation}
\label{sec:reflow}

After training a basic flow matching model $v_1$, you can straighten the
trajectories further via the reflow procedure.

\subsection{Generating Coupled Pairs}

\begin{tcolorbox}[colback=codebg, colframe=codeframe, title=\textbf{Step 1: Generate coupled pairs using $v_1$}]
\begin{lstlisting}
# Generate noise samples
x_0 = torch.randn(N, dim)

# Solve ODE forward from t=0 to t=1 using learned v_1
x = x_0.clone()
dt = 1.0 / n_steps
for i in range(n_steps):
    t = i * dt
    x = x + v_1(x, t) * dt
z_1 = x  # coupled output

# Store (x_0, z_1) pairs -- these are now correlated
\end{lstlisting}
\end{tcolorbox}

\subsection{Retraining on Coupled Pairs}

\begin{tcolorbox}[colback=codebg, colframe=codeframe, title=\textbf{Step 2: Retrain $v_2$ on coupled pairs}]
\begin{lstlisting}
for each batch:
    # Use COUPLED pairs instead of independent samples
    x_0, z_1 = sample_coupled_pairs()
    t = torch.rand(batch_size, 1)
    x_t = (1 - t) * x_0 + t * z_1
    v_target = z_1 - x_0
    v_pred = model_v2(x_t, t)
    loss = MSE(v_pred, v_target)
    ...
\end{lstlisting}
\end{tcolorbox}

\subsection{One-Step Generation After Reflow}

After sufficient reflow rounds, the trajectories are approximately straight,
enabling:
\begin{equation}
  x_{\mathrm{generated}} = x_0 + v_k(x_0, 0),
  \qquad x_0 \sim \N(0, \Id)
\end{equation}

This replaces the entire multi-step ODE integration with a single function
evaluation.

% ══════════════════════════════════════════════════════════════
\section{Reading Order and Resources}
\label{sec:resources}

\subsection{Recommended Reading Order}

\begin{enumerate}
  \item \textbf{Start here:}
    \texttt{THEORY/GUIDE\_Flow\_Matching\_for\_ToyDiffusion.md} \\
    Practical, code-focused, maps directly to ToyDiffusion.

  \item \textbf{Intuition:}
    Cambridge MLG Blog \\
    \url{https://mlg.eng.cam.ac.uk/blog/2024/01/20/flow-matching.html} \\
    Builds from CNFs $\to$ FM $\to$ CFM $\to$ OT paths.

  \item \textbf{Interactive visuals:}
    Diffusion Meets Flow Matching \\
    \url{https://diffusionflow.github.io/} \\
    Mathematical equivalence shown visually.

  \item \textbf{Paper 1:} Flow Matching (Lipman et al., 2022) \\
    \url{https://arxiv.org/abs/2210.02747} \\
    Read Sections 1--3 for the math.

  \item \textbf{Paper 2:} Rectified Flow (Liu et al., 2022) \\
    \url{https://arxiv.org/abs/2209.03003} \\
    Read Sections 1--3 for the reflow procedure.

  \item \textbf{Implementation companion:} Flow Matching Guide and Code
    (Lipman et al., 2024) \\
    \url{https://arxiv.org/abs/2412.06264} \\
    Comprehensive tutorial with PyTorch code.

  \item \textbf{MIT Course:} Flow Matching and Diffusion Models (2026) \\
    \url{https://diffusion.csail.mit.edu/docs/lecture-notes.pdf}
\end{enumerate}

\subsection{Code Repositories}

\begin{center}
\begin{tabular}{@{} l l l @{}}
\toprule
\textbf{Repo} & \textbf{What} & \textbf{Link} \\
\midrule
RectifiedFlow & Official PyTorch impl.
  & \url{https://github.com/gnobitab/RectifiedFlow} \\
ToyDiffusion & Your original DDPM project
  & \url{https://github.com/ThiagoLira/ToyDiffusion} \\
\bottomrule
\end{tabular}
\end{center}

% ══════════════════════════════════════════════════════════════
\newpage
\section*{Flow Matching Cheat Sheet}
\addcontentsline{toc}{section}{Flow Matching Cheat Sheet}

\begin{tcolorbox}[
  colback=fmgreen!5, colframe=fmgreen!80!black,
  fonttitle=\bfseries\large,
  title=Quick Reference Card,
  rounded corners, boxrule=1pt
]

\textbf{Training:}
\begin{align*}
  x_1 &\sim p_{\mathrm{data}} && \text{sample data} \\
  x_0 &\sim \N(0, \Id) && \text{sample noise} \\
  t   &\sim U(0, 1) && \text{sample time} \\
  x_t &= (1-t)\,x_0 + t\,x_1 && \text{interpolate} \\
  v   &= x_1 - x_0 && \text{target velocity} \\
  \mathcal{L} &= \norm{\vtheta(x_t, t) - v}^2 && \text{MSE loss}
\end{align*}

\textbf{Sampling} (Euler, $N$ steps):
\begin{align*}
  x &\sim \N(0, \Id), \qquad \Delta t = 1/N \\
  \text{for } i &= 0, \ldots, N{-}1: \\
  &\quad t = i \cdot \Delta t \\
  &\quad x \leftarrow x + \vtheta(x, t) \cdot \Delta t \\
  &\text{return } x
\end{align*}

\textbf{What the network predicts:}
\[
  \text{Input: } (x_t, t) \quad \longrightarrow \quad
  \text{Output: } v \in \R^d \text{ (velocity, same dim as $x_t$)}
\]

\medskip
\begin{center}
\large\bfseries
That's it.  No $\beta$s.  No $\alpha$s.  No posterior.  Just velocity.
\end{center}

\end{tcolorbox}

% ══════════════════════════════════════════════════════════════
\section*{References}
\addcontentsline{toc}{section}{References}

\begin{enumerate}[label={[\arabic*]}]
  \item J.~Ho, A.~Jain, P.~Abbeel.
    \textit{Denoising Diffusion Probabilistic Models}.
    NeurIPS 2020.
    \url{https://arxiv.org/abs/2006.11239}

  \item Y.~Lipman, R.~T.~Q.~Chen, H.~Ben-Hamu, M.~Nickel, M.~Le.
    \textit{Flow Matching for Generative Modeling}.
    ICLR 2023.
    \url{https://arxiv.org/abs/2210.02747}

  \item X.~Liu, C.~Gong, Q.~Liu.
    \textit{Flow Straight and Fast: Learning to Generate and Transfer
    Data with Rectified Flow}.
    ICLR 2023 (Spotlight).
    \url{https://arxiv.org/abs/2209.03003}

  \item Y.~Lipman, M.~Havasi, P.~Holderrieth, N.~Shaul, M.~Le,
    B.~Karrer, R.~T.~Q.~Chen, D.~Lopez-Paz, H.~Ben-Hamu, I.~Gat.
    \textit{Flow Matching Guide and Code}.
    2024.
    \url{https://arxiv.org/abs/2412.06264}

  \item A.~Vaswani, N.~Shazeer, N.~Parmar, et al.
    \textit{Attention Is All You Need}.
    NeurIPS 2017.

  \item J.~Song, C.~Meng, S.~Ermon.
    \textit{Denoising Diffusion Implicit Models (DDIM)}.
    ICLR 2021.
    \url{https://arxiv.org/abs/2010.02502}
\end{enumerate}

\end{document}
